% =============================================================================
%  سجل متابعة المخبر — مادة المعلوماتية
%  السنة الدراسية: 2026 / 2027
% =============================================================================
%  البناء (من هذا المجلد):
%    xelatex -interaction=nonstopmode -halt-on-error -output-directory=build lab-register.tex
%  (مرّتان أو ثلاث مرّات لأنّ التذييل يستعمل \pageref{LastPage})
% =============================================================================

\documentclass[11pt, a4paper]{article}

% ── المشترك (مسار shared نسبي لمجلد documents/) ──────────────────────────────
\def\shareddir{../shared/}
\input{"../shared/common.tex"}

% ── معلومات الملف ────────────────────────────────────────────────────────────
\hypersetup{
  pdfauthor={لفقير عبدالجليل},
  pdftitle={سجل متابعة المخبر},
}

% ══════════════════════════════════════════════════════════════════════════════
%  ماكروهات سجل المخبر
% ══════════════════════════════════════════════════════════════════════════════

% ── عنوان قسم: صفحة جديدة + ترويسة رسمية + عنوان ────────────────────────────
\newcommand{\regsection}[1]{%
  \clearpage
  \officialheader
  \vspace{2pt}%
  \begin{center}{\LARGE\bfseries\color{primary} #1}\end{center}
  \vspace{8pt}%
}

% ── صفوف مكرّرة للجداول الممتدة ─────────────────────────────────────────────
\newcommand{\devrow}{& & & & & & \\[0.30cm] \hline}
\newcommand{\devrows}[1]{\ifnum#1>0 \devrow \devrows{\number\numexpr#1-1\relax}\fi}

\newcommand{\logline}{& & & & & & \\[0.38cm] \hline}
\newcommand{\loglines}[1]{\ifnum#1>0 \logline \loglines{\number\numexpr#1-1\relax}\fi}

\newcommand{\seatrow}{& & \\[0.02cm] \hline}
\newcommand{\seatrows}[1]{\ifnum#1>0 \seatrow \seatrows{\number\numexpr#1-1\relax}\fi}

% ── جدول الجرد ───────────────────────────────────────────────────────────────
\newcommand{\devinventory}{%
  \begingroup
  \gridsetup
  \renewcommand{\arraystretch}{1.2}%
  \setlength{\tabcolsep}{2pt}%
  \begin{longtable}{@{}|c|p{1.8cm}|p{4.2cm}|p{2.7cm}|c|p{2.3cm}|p{2.7cm}|@{}}
    \hline
    \rowcolor{primary}
    \thead[0.8cm]{الرقم} & \thead[1.8cm]{رقم الجرد} & \thead[4.2cm]{تسمية العتاد} &
    \thead[2.7cm]{الطراز / المواصفات} & \thead[0.8cm]{العدد} &
    \thead[2.3cm]{الحالة} & \thead[2.7cm]{ملاحظات} \\
    \hline
    \endfirsthead
    \hline
    \rowcolor{primary}
    \thead[0.8cm]{الرقم} & \thead[1.8cm]{رقم الجرد} & \thead[4.2cm]{تسمية العتاد} &
    \thead[2.7cm]{الطراز / المواصفات} & \thead[0.8cm]{العدد} &
    \thead[2.3cm]{الحالة} & \thead[2.7cm]{ملاحظات} \\
    \hline
    \endhead
    \devrows{20}
  \end{longtable}
  \endgroup
}

% ── جدول سجل المتابعة ────────────────────────────────────────────────────────
\newcommand{\logtable}{%
  \begingroup
  \gridsetup
  \setlength{\tabcolsep}{3pt}%
  \begin{longtable}{@{}|p{1.9cm}|p{1.5cm}|p{2.4cm}|p{3.0cm}|p{2.4cm}|p{1.8cm}|p{2.0cm}|@{}}
    \hline
    \rowcolor{primary}
    \thead{التاريخ} & \thead{الحصة} & \thead{القسم} & \thead{النشاط} &
    \thead{العتاد المستعمل} & \thead{الحالة} & \thead{الملاحظات} \\
    \hline
    \endfirsthead
    \hline
    \rowcolor{primary}
    \thead{التاريخ} & \thead{الحصة} & \thead{القسم} & \thead{النشاط} &
    \thead{العتاد المستعمل} & \thead{الحالة} & \thead{الملاحظات} \\
    \hline
    \endhead
    \loglines{72}
  \end{longtable}
  \endgroup
}

% ── جدول الرزنامة (الأيام × الحصص) ──────────────────────────────────────────
\newcommand{\schedrow}[1]{%
  \cellcenter{2.0cm}{\textbf{\color{primary}#1}} & & & & & & & \\[1.35cm] \hline}
\newcommand{\schedule}{%
  \begingroup
  \gridsetup
  \setlength{\tabcolsep}{3pt}%
  \begin{tabularx}{\textwidth}{@{}|c|X|X|X|X|X|X|X|@{}}
    \hline
    \rowcolor{primary}
    \thead[2.0cm]{اليوم} & \thead{الحصة 1} & \thead{الحصة 2} & \thead{الحصة 3} &
    \thead{الحصة 4} & \thead{الحصة 5} & \thead{الحصة 6} & \thead{الحصة 7} \\
    \hline
    \schedrow{الأحد}
    \schedrow{الاثنين}
    \schedrow{الثلاثاء}
    \schedrow{الأربعاء}
    \schedrow{الخميس}
  \end{tabularx}
  \endgroup
}

% ── قائمة جلوس قسم واحد ─────────────────────────────────────────────────────
\newcommand{\seatinglist}[1]{%
  \clearpage
  \officialheader
  \vspace{2pt}%
  \begin{center}
    {\large\bfseries\color{primary} قائمة جلوس التلاميذ
     — \textcolor{secondary}{#1}}
  \end{center}
  \vspace{4pt}
  \begingroup
  \gridsetup
  \renewcommand{\arraystretch}{1.05}%
  \setlength{\tabcolsep}{4pt}%
  \begin{longtable}{@{}|c|p{7.6cm}|p{2.6cm}|p{4.0cm}|@{}}
    \hline
    \rowcolor{primary}
    \thead[1.0cm]{الرقم} & \thead{اللقب والاسم} & \thead[2.6cm]{رقم المحطة} &
    \thead{الملاحظات} \\
    \hline
    \endfirsthead
    \hline
    \rowcolor{primary}
    \thead[1.0cm]{الرقم} & \thead{اللقب والاسم} & \thead[2.6cm]{رقم المحطة} &
    \thead{الملاحظات} \\
    \hline
    \endhead
    \seatrows{30}
  \end{longtable}
  \endgroup
}

% تعطيل التظليل المتناوب داخل الجداول الممتدة (تفادي تداخله مع longtable)
\rowcolors{1}{}{}

% ══════════════════════════════════════════════════════════════════════════════
\begin{document}

% ═════════════════════════════════════════════════════════════════════════════
%  الغلاف
% ═════════════════════════════════════════════════════════════════════════════
\officialheader

\vspace*{\fill}
\begin{center}
  {\Huge\bfseries\color{primary} سجل متابعة المخبر}\\[12pt]
  {\LARGE\bfseries\color{secondary} مخبر الإعلام الآلي}\\[8pt]
  {\large مادة المعلوماتية}\\[4pt]
  {\normalsize\color{primary} السنة الدراسية 2026 / 2027}
\end{center}
\vspace*{\fill}

% ═════════════════════════════════════════════════════════════════════════════
%  1. جرد الأجهزة
% ═════════════════════════════════════════════════════════════════════════════
\regsection{جرد الأجهزة}

{\small يُدوَّن في الجدول الآتي كل عتاد المخبر (وحدات مركزية، شاشات، لوحات مفاتيح،
فأرات، طابعات، مجهّزات الشبكة،\,\ldots) مع عدد كل صنف وحالته.}

\vspace{6pt}
\devinventory

\vspace{6pt}
\begin{rememberbox}[ملاحظة]
  يُراجَع هذا الجرد في بداية كل سنة دراسية، وتُدوَّن فيه كل عمليات الاقتناء
  والإصلاح والإخراج، مع الإشارة إلى العتاد المعطّل ورقم الوصل عند الاقتضاء.
\end{rememberbox}

% ═════════════════════════════════════════════════════════════════════════════
%  2. القانون الداخلي للمخبر
% ═════════════════════════════════════════════════════════════════════════════
\regsection{القانون الداخلي للمخبر}

{\small يلتزم كل مستعمل للمخبر (أستاذًا كان أو تلميذًا) بما يأتي:}

\vspace{6pt}
{\small
\begin{enumerate}[label=\textcolor{primary}{\bfseries\arabic*.},
    labelwidth=1.6em, left=0.4em .. 2.1em, itemsep=7pt, parsep=0pt, topsep=2pt]
  \item الالتزام بالحضور إلى المخبر في التوقيت المحدّد، والدخول والخروج بهدوء ونظام.
  \item الجلوس في المكان المخصّص لكل تلميذ وفق قائمة الجلوس، وعدم تغييره إلا بإذن الأستاذ.
  \item تشغيل الأجهزة وإطفاؤها وفق تعليمات الأستاذ فقط.
  \item عدم لمس الأسلاك والوصلات الكهربائية أو فتح وحدات الحاسوب أو نزع مكوّناتها.
  \item عدم إدخال أو استعمال أي وسائط تخزين (USB، أقراص مدمجة) دون إذن مسبق من الأستاذ.
  \item عدم تنصيب أو حذف أي برنامج، وعدم تغيير إعدادات النظام أو كلمات المرور.
  \item المحافظة على العتاد والطاولات والكراسي، والإبلاغ الفوري عن أي عطب أو خلل.
  \item يُمنع الأكل والشرب داخل المخبر.
  \item يُمنع استعمال الهاتف أو أي جهاز إلكتروني إلا في إطار النشاط التعليمي وبإذن الأستاذ.
  \item استعمال الشبكة والإنترنت في الأغراض الدراسية فقط.
  \item عدم ربط أي جهاز شخصي بالمخبر أو تغيير مكان العتاد دون إذن الأستاذ.
  \item الحفاظ على نظافة المخبر وترتيبه، وترك المكان نظيفًا مرتبًا عند الخروج.
  \item احترام الأستاذ والزملاء، والالتزام بالهدوء أثناء العمل.
  \item تعويض كل عتاد يتلف نتيجة الإهمال أو الاستعمال غير السليم.
\end{enumerate}\par}

% ═════════════════════════════════════════════════════════════════════════════
%  3. سجل متابعة المخبر
% ═════════════════════════════════════════════════════════════════════════════
\regsection{سجل متابعة المخبر}

{\small تُدوَّن في هذا السجل كل حصة استُعمل فيها المخبر: التاريخ، الحصة، القسم،
النشاط المنجَز، العتاد المستعمل، وحالة المخبر عند الخروج.}

\vspace{6pt}
\logtable

% ═════════════════════════════════════════════════════════════════════════════
%  4. رزنامة استغلال المخبر
% ═════════════════════════════════════════════════════════════════════════════
\regsection{رزنامة استغلال المخبر}

{\small تُعبَّأ الخلايا باسم القسم (أو الأستاذ) المستفيد من المخبر في كل حصة من حصص الأسبوع.}

\vspace{8pt}
\schedule

% ═════════════════════════════════════════════════════════════════════════════
%  5. قوائم جلوس التلاميذ
% ═════════════════════════════════════════════════════════════════════════════
\seatinglist{السنة الأولى جذع مشترك علوم وتكنولوجيا — الفوج 1}
\seatinglist{السنة الأولى جذع مشترك علوم وتكنولوجيا — الفوج 2}
\seatinglist{السنة الأولى جذع مشترك علوم وتكنولوجيا — الفوج 3}
\seatinglist{السنة الأولى جذع مشترك علوم وتكنولوجيا — الفوج 4}
\seatinglist{السنة الأولى جذع مشترك علوم وتكنولوجيا — الفوج 5}
\seatinglist{السنة الأولى جذع مشترك آداب وفلسفة — الفوج 1}
\seatinglist{السنة الأولى جذع مشترك آداب وفلسفة — الفوج 2}
\seatinglist{السنة الأولى جذع مشترك آداب وفلسفة — الفوج 3}

\end{document}
